\documentclass[11pt]{article}

\usepackage[margin=1in]{geometry}
\usepackage[T1]{fontenc}
\usepackage[utf8]{inputenc}
\usepackage{lmodern}
\usepackage{booktabs}
\usepackage{enumitem}
\usepackage{hyperref}
\usepackage{longtable}

\hypersetup{
  colorlinks=true,
  linkcolor=blue!50!black,
  urlcolor=blue!50!black,
  citecolor=blue!50!black,
  pdftitle={Prior-Art Pass Report for p-vs-np}
}

\title{Prior-Art Pass Report for \texttt{p-vs-np}}
\author{Neuro-symbolic AI Group\\Idiap Research Institute}
\date{March 17, 2026}

\begin{document}
\maketitle

\begin{abstract}
This report records a targeted prior-art pass over the current blocked proof frontier. It imports decisive external theorem-level references into the problem-scoped InformalMathLib and explains what those references change about the current proof state.
\end{abstract}

\section{Snapshot}
\begin{itemize}[leftmargin=*,itemsep=0.2em]
\item Imported references: 8
\item Blocked queue items: 24
\item Covered blocked queue items: 24
\item Uncovered blocked queue items: 
\end{itemize}

\section{Search Justification}
The local queue is exhausted through route conjectures, abstraction repairs, and repair lemmas, so browsing is justified only for named missing lemmas that materially affect novelty or barrier judgments.

\section{Key Findings}
\begin{itemize}[leftmargin=*,itemsep=0.2em]
\item The imported references ground most target-side theorems and barriers, but none closes the missing carrier-realization or formula-size transport lemmas.
\item The PIT route sharpens from a vague semantic-to-algorithmic gap into an explicitness gap: known results already connect explicit hitting sets to black-box PIT and PIT derandomization to lower bounds.
\item The Cook-Reckhow route is not discovering a new target theorem; the real gap is the bridge from the local definitional object to the existing NP=coNP characterization.
\item The MCSP local-reduction route is now barrier-aware because prior art already shows non-NP-hardness under local reductions and strong consequences for harder reductions.
\item The general magnification route is literature-backed, but the unresolved step remains object-class alignment between the MCSP-specific carrier and the sparse-problem or distinguisher framework.
\end{itemize}

\section{Imported References}
\begin{longtable}{p{0.13\linewidth}p{0.1\linewidth}p{0.42\linewidth}p{0.25\linewidth}}
\toprule
ID & Type & Title & Obstruction classes \\
\midrule
mathlibref.p-vs-np.001 & theorem & Hardness Magnification for Natural Problems & carrier-realization, mcsp-magnification \\
mathlibref.p-vs-np.002 & theorem & Hardness Magnification Near State-of-the-Art Lower Bounds & carrier-realization, formula-size-transport, mcsp-magnification \\
mathlibref.p-vs-np.003 & theorem & MCSP Is Not NP-Hard Under Local Reductions & local-reduction-barrier \\
mathlibref.p-vs-np.004 & theorem & Hitting-Set Generators Yield Black-Box PIT & pit-upgrade, explicitness-gap \\
mathlibref.p-vs-np.005 & theorem & Derandomizing Polynomial Identity Tests Means Proving Circuit Lower Bounds & pit-upgrade, target-grounding \\
mathlibref.p-vs-np.008 & theorem & Many Natural Problems Are NP-Complete & schema-to-examples, target-grounding \\
mathlibref.p-vs-np.006 & theorem & Polynomially Bounded Proof Systems Characterize NP Equals coNP & characterization-bridge, target-grounding \\
mathlibref.p-vs-np.007 & theorem & A Simple and General Hardness Magnification Theorem & target-alignment, general-magnification \\
\bottomrule
\end{longtable}

\section{Obstruction Map}
\subsection*{carrier-realization}
\textbf{References.} mathlibref.p-vs-np.001, mathlibref.p-vs-np.002

\textbf{Queue ids.} proofq.p-vs-np.001, proofq.p-vs-np.002, proofq.p-vs-np.005, proofq.p-vs-np.009, proofq.p-vs-np.010, proofq.p-vs-np.013, proofq.p-vs-np.018, proofq.p-vs-np.020, proofq.p-vs-np.021

\begin{itemize}[leftmargin=*,itemsep=0.2em]
\item The leading MCSP route is no longer novel at the level of magnification-as-such; the open step is the typed carrier or realization theorem that would instantiate the route inside the current proof state.
\item The bridge.001 target is externally grounded, so the remaining novelty claim shifts to how the current framework realizes a usable MCSP carrier rather than to the target theorem itself.
\end{itemize}
\subsection*{characterization-bridge}
\textbf{References.} mathlibref.p-vs-np.006

\textbf{Queue ids.} proofq.p-vs-np.006, proofq.p-vs-np.014, proofq.p-vs-np.024

\begin{itemize}[leftmargin=*,itemsep=0.2em]
\item The Cook-Reckhow route is now clearly a bridge problem, not a target-theorem problem: the open step is the semantic bridge from the local definitional object to this existing characterization theorem.
\end{itemize}
\subsection*{explicitness-gap}
\textbf{References.} mathlibref.p-vs-np.004

\textbf{Queue ids.} proofq.p-vs-np.016, proofq.p-vs-np.017

\begin{itemize}[leftmargin=*,itemsep=0.2em]
\item The PIT route is not blocked by the total absence of a hitting-set-to-PIT connection; it is blocked by the lack of an explicit generator or certificate object for the right class.
\end{itemize}
\subsection*{formula-size-transport}
\textbf{References.} mathlibref.p-vs-np.002

\textbf{Queue ids.} proofq.p-vs-np.001, proofq.p-vs-np.005, proofq.p-vs-np.009, proofq.p-vs-np.010, proofq.p-vs-np.013, proofq.p-vs-np.018, proofq.p-vs-np.020, proofq.p-vs-np.021

\begin{itemize}[leftmargin=*,itemsep=0.2em]
\item The bridge.001 target is externally grounded, so the remaining novelty claim shifts to how the current framework realizes a usable MCSP carrier rather than to the target theorem itself.
\end{itemize}
\subsection*{general-magnification}
\textbf{References.} mathlibref.p-vs-np.007

\textbf{Queue ids.} proofq.p-vs-np.003, proofq.p-vs-np.011, proofq.p-vs-np.023

\begin{itemize}[leftmargin=*,itemsep=0.2em]
\item The general magnification route is literature-backed, but its live gap is an alignment theorem from the MCSP obstruction family to the sparse-problem or distinguisher object class required by the general theorem.
\end{itemize}
\subsection*{local-reduction-barrier}
\textbf{References.} mathlibref.p-vs-np.003

\textbf{Queue ids.} proofq.p-vs-np.004, proofq.p-vs-np.012, proofq.p-vs-np.019

\begin{itemize}[leftmargin=*,itemsep=0.2em]
\item The local-reduction route is now explicitly barrier-aware: the missing gadget lemma must overcome known non-hardness and consequence results, not just fill an arbitrary construction gap.
\end{itemize}
\subsection*{mcsp-magnification}
\textbf{References.} mathlibref.p-vs-np.001, mathlibref.p-vs-np.002

\textbf{Queue ids.} proofq.p-vs-np.001, proofq.p-vs-np.002, proofq.p-vs-np.005, proofq.p-vs-np.009, proofq.p-vs-np.010, proofq.p-vs-np.013, proofq.p-vs-np.018, proofq.p-vs-np.020, proofq.p-vs-np.021

\begin{itemize}[leftmargin=*,itemsep=0.2em]
\item The leading MCSP route is no longer novel at the level of magnification-as-such; the open step is the typed carrier or realization theorem that would instantiate the route inside the current proof state.
\item The bridge.001 target is externally grounded, so the remaining novelty claim shifts to how the current framework realizes a usable MCSP carrier rather than to the target theorem itself.
\end{itemize}
\subsection*{pit-upgrade}
\textbf{References.} mathlibref.p-vs-np.004, mathlibref.p-vs-np.005

\textbf{Queue ids.} proofq.p-vs-np.008, proofq.p-vs-np.016, proofq.p-vs-np.017

\begin{itemize}[leftmargin=*,itemsep=0.2em]
\item The PIT route is not blocked by the total absence of a hitting-set-to-PIT connection; it is blocked by the lack of an explicit generator or certificate object for the right class.
\item The PIT route's target theorem is prior art; the remaining open step is the current framework's inability to produce the required derandomization premise from its local PIT carrier.
\end{itemize}
\subsection*{schema-to-examples}
\textbf{References.} mathlibref.p-vs-np.008

\textbf{Queue ids.} proofq.p-vs-np.007, proofq.p-vs-np.015, proofq.p-vs-np.022

\begin{itemize}[leftmargin=*,itemsep=0.2em]
\item The NP-complete examples route is now target-grounded: the open step is constructing or transferring explicit examples from the local abstraction, not discovering that natural NP-complete problems exist.
\end{itemize}
\subsection*{target-alignment}
\textbf{References.} mathlibref.p-vs-np.007

\textbf{Queue ids.} proofq.p-vs-np.003, proofq.p-vs-np.011, proofq.p-vs-np.023

\begin{itemize}[leftmargin=*,itemsep=0.2em]
\item The general magnification route is literature-backed, but its live gap is an alignment theorem from the MCSP obstruction family to the sparse-problem or distinguisher object class required by the general theorem.
\end{itemize}
\subsection*{target-grounding}
\textbf{References.} mathlibref.p-vs-np.005, mathlibref.p-vs-np.006, mathlibref.p-vs-np.008

\textbf{Queue ids.} proofq.p-vs-np.006, proofq.p-vs-np.007, proofq.p-vs-np.008, proofq.p-vs-np.014, proofq.p-vs-np.015, proofq.p-vs-np.016, proofq.p-vs-np.017, proofq.p-vs-np.022, proofq.p-vs-np.024

\begin{itemize}[leftmargin=*,itemsep=0.2em]
\item The PIT route's target theorem is prior art; the remaining open step is the current framework's inability to produce the required derandomization premise from its local PIT carrier.
\item The NP-complete examples route is now target-grounded: the open step is constructing or transferring explicit examples from the local abstraction, not discovering that natural NP-complete problems exist.
\item The Cook-Reckhow route is now clearly a bridge problem, not a target-theorem problem: the open step is the semantic bridge from the local definitional object to this existing characterization theorem.
\end{itemize}

\section{Still Uncovered}
\begin{itemize}[leftmargin=*,itemsep=0.2em]
\item All blocked queue items have at least one linked imported reference.
\end{itemize}

\end{document}
